\exercisegroup

\begin{enumerate}
\def\labelenumi{\arabic{enumi}.}
\item
  设 I 是一个有向集， \((\mathbf{J}_{\lambda})_{\lambda \in \mathbf{L}}\) 是由有向集 L 索引的 I 的子集族，并满足：(i) 对每个 \(\lambda \in \mathbf{L}\) , \(\mathbf{J}_{\lambda}\) 关于诱导序是有向的；(ii) 关系 \(\lambda \leqslant \mu\) 蕴含 \(\mathbf{J}_{\lambda} \subset \mathbf{J}_{\mu}\) ；(iii) I 是族 \((\mathbf{J}_{\lambda})\) 的并。设 \((\mathbf{E}_{\alpha}, f_{\alpha\beta})\) 是关于 I 的集合逆系统，E 为其逆极限；对每个 \(\lambda \in \mathbf{L}\) ，令 \(\mathbf{F}_{\lambda}\) 为把 \((\mathbf{E}_{\alpha}, f_{\alpha\beta})\) 的指标集限制到 \(\mathbf{J}_{\lambda}\) 后所得系统的逆极限。对 \(\lambda \leqslant \mu\) ，令 \(g_{\lambda \mu}\) 为从 \(\mathbf{F}_{\mu}\) 到 \(\mathbf{F}_{\lambda}\) 的典范映射（第1号）。证明 \((\mathbf{F}_{\lambda}, g_{\lambda \mu})\) 是关于 L 的集合逆系统，并定义一个从 \(\mathbf{F} = \lim \mathbf{F}_{\lambda}\) 到 E 的典范双射。
\item
  让 \((\mathbf{E}_{\alpha}, f_{\alpha\beta})\) 是相对于一个有向指标集的集合逆系统，设 \(E = \lim E_{\alpha}\) ，并且让 \(f_{\alpha}: E \to E_{\alpha}\) 为对每个 \(\alpha\) 而言的典范映射。证明：如果所有 \(f_{\alpha\beta}\) 都是单射的，那么 \(f_{\alpha}\) 是单射的。
\item
  设 \((\mathbf{E}_{\alpha}, f_{\alpha\beta})\) 和 \((\mathbf{F}_{\alpha}, g_{\alpha\beta})\) 是关于同一指标集 I 的两个集合逆系统。对每个 \(\alpha \in I\) ，令 \(u_{\alpha}\) 为从 \(E_{\alpha}\) 到 \(F_{\alpha}\) 的映射，并使各 \(u_{\alpha}\) 构成一个映射逆系统。设 \(G_{\alpha} \subset E_{\alpha} \times F_{\alpha}\) 是 \(u_{\alpha}\) 的图。证明 \((\mathbf{G}_{\alpha})\) 是各 \(E_{\alpha} \times F_{\alpha}\) 的一个子集逆系统，并且其逆极限可典范地等同于 \(u = \lim u_{\alpha}\) 的图。
\item
  设 I 是一个没有最大元素的非空有向集，并设 F 为所有满足以下性质的序列的集合： \(x = (\alpha_1, \alpha_2, \ldots, \alpha_{2n-1}, \alpha_{2n})\) 由偶数个元素组成 \(\geqslant 2\) 且元素取自 I：(i) \(\alpha_{2i-1} < \alpha_{2i}\) 对于 \(1 \leqslant i \leqslant n\) ；(ii) \(\alpha_{2i-1} \leqslant \alpha_{2j-1}\) 对于 \(1 \leqslant j < i \leqslant n\) 。集合 F 非空。令 \(r(x) = \alpha_{2n-1}\) , \(s(x) = \alpha_{2n}\) 。整数 \(n\) 称为 \(x\) 的长度.
\end{enumerate}

\begin{enumerate}
\def\labelenumi{(\alph{enumi})}
\tightlist
\item
  对于每一个 \(\alpha \in I\) ，让 \(E_{\alpha}\) 是由所有这样的 \(x \in F\) （满足 \(r(x) = \alpha\) ）所组成的集合。那么 \(E_{\alpha}\) 不为空。对于 \(\alpha \leqslant \beta\) （在 \(I\) 中），我们定义一个映射 \(f_{\alpha \beta}\) （从 \(E_{\beta}\) 到由 \(I\) 的元素组成的所有有限序列的集合中），如下：如果
\end{enumerate}

\[
x = \left(\alpha_ {1}, \dots , \alpha_ {2 n - 1}, \alpha_ {2 n}\right) \in \mathrm{E} _ {\beta},
\]

让 \(j\) 是使得 \(\alpha \leqslant \alpha_{2j - 1}\) 的最小指标；然后

\[
f _ {\alpha \beta} (x) = (\alpha_ {1}, \dots , \alpha_ {2 j - 2}, \alpha , \alpha_ {2 j}).
\]

证明 \(f_{\alpha \beta}(\mathrm{E}_{\beta}) = \mathrm{E}_{\alpha}\) 以及 \((\mathrm{E}_{\alpha}, f_{\alpha \beta})\) 是相对于 \(\mathbf{I}\) 的集合逆系统.

\begin{enumerate}
\def\labelenumi{(\alph{enumi})}
\setcounter{enumi}{1}
\item
  证明如果 \(x_{\alpha} \in \mathbf{E}_{\alpha}\) 并且 \(x_{\beta} \in \mathbf{E}_{\beta}\) 使得存在一个指标 \(\gamma\) 满足 \(\gamma \geqslant \alpha\) 并且 \(\gamma \geqslant \beta\) ，以及一个元素 \(x_{\gamma} \in \mathbf{E}_{\gamma}\) 满足 \(x_{\alpha} = f_{\alpha \gamma}(x_{\gamma})\) 并且 \(x_{\beta} = f_{\beta \gamma}(x_{\gamma})\) ，那么，在 \(x_{\alpha}\) 与 \(x_{\beta}\) 具有相同长度的条件下，我们有 \(s(x_{\alpha}) = s(x_{\beta})\) .
\item
  从（b）推出，如果 \(\mathbf{E} = \lim \mathbf{E}_{\alpha}\) 非空，且若 \(y = (x_{\alpha}) \in \mathbf{E}\) ，则元素 \(s(x_{\alpha})\) 所组成的集合是可数的且在 \(\mathbf{I}\) 中共尾.
\item
  设 I 为不可数集合 A 的所有有限子集构成的集合，按包含关系排序。证明 I 没有可数的共尾子集，并由此从 (c) 推出一个集合逆系统的例子 \((\mathbf{E}_{\alpha}, f_{\alpha\beta})\) ，在其中 \(E_{\alpha}\) 是非空的，且 \(f_{\alpha\beta}\) 是满射的，但对其有 \(E = \lim E_{\alpha} = \phi\) .
\item
  从(d)推导出一个映射逆系统的例子 \(u_{\alpha}: \mathbf{E}_{\alpha} \to \mathbf{E}_{\alpha}'\) ，使得每个 \(u_{\alpha}\) 是满射但 \(\lim u_{\alpha}\) 不是满射（令每个 \(\mathbf{E}_{\alpha}'\) 由单个元素组成）。
\end{enumerate}

\begin{enumerate}
\def\labelenumi{\arabic{enumi}.}
\setcounter{enumi}{4}
\tightlist
\item
  设 I 是一个有向集，并令 \((\mathbf{E}_{\alpha})_{\alpha \in \mathbf{I}}\) 是一族格，使得每个 \(\mathbf{E}_{\alpha}\) ，赋予相反的序后，是诺特集（§ 6，第 5 号）。对于每一对指标 \((\alpha, \beta)\) （ I 中满足 \(\alpha \leqslant \beta\) 者），让
\end{enumerate}

\[
f _ {\alpha \beta}: \mathbf {E} _ {\beta} \rightarrow \mathbf {E} _ {\alpha}
\]

是一个递增映射，并假设 \((\mathbf{E}_{\alpha}, f_{\alpha\beta})\) 是相对于 I 的集合的逆向系统。对每个 \(\alpha \in I\) 让 \(G_{\alpha}\) 是 \(E_{\alpha}\) 的非空子集，使得 (i) \(G_{\alpha}\) 中任何两个不同的元素都不是可比较的，(ii) \(f_{\alpha\beta}(G_{\beta}) = G_{\alpha}\) 每当 \(\alpha \leqslant \beta\) ， (iii) 对每个 \(\alpha \leqslant \beta\) 以及每个 \(x_{\alpha} \in G_{\alpha}\) , \(f_{\alpha\beta}(x_{\alpha})\) 有最大元素 \(M_{\alpha\beta}(x_{\alpha})\) （在 \(E_{\beta}\) 中），(iv) 每当 \(\alpha \leqslant \beta\) ，如果 \(h_{\beta} \in E_{\beta}\) 使得存在 \(y_{\beta} \in G_{\beta}\) 使得 \(y_{\beta} \leqslant h_{\beta}\) ，那么对于每个 \(x_{\alpha} \in G_{\alpha}\) （满足 \(x_{\alpha} \leqslant f_{\alpha\beta}(h_{\beta})\) ），存在 \(x_{\beta} \in G_{\beta}\) 使得 \(x_{\beta} \leqslant h_{\beta}\) 并且 \(x_{\alpha} = f_{\alpha\beta}(x_{\beta})\) 。在这些条件下，子集逆系统 \((\mathbf{G}_{\alpha})\) 的逆极限非空。证明如下：

\begin{enumerate}
\def\labelenumi{(\alph{enumi})}
\item
  让 \(J\) 是 \(I\) 的一个有限子集。一个族 \((x_{\alpha})_{\alpha \in J}\) ，其中 \(x_{\alpha} \in G_{\alpha}\) 对于所有 \(\alpha \in J\) ，如果它满足以下两个条件，则称其为一致的：
\item
  （i）若 \(\alpha \in J\) , \(\beta \in J\) , \(\alpha \leqslant \beta\) ，那么 \(x_{\alpha} = f_{\alpha \beta}(x_{\beta})\) ； (ii) 对于每个上界 \(\gamma\) （ \(J\) 在 \(I\) 中的上界），存在 \(x_{\gamma} \in G_{\gamma}\) 使得 \(x_{\alpha} = f_{\alpha \gamma}(x_{\gamma})\) 对于所有 \(\alpha \in J\) 。证明，对于每个上界 \(\gamma\) （ \(J\) 在 \(I\) 中的上界），集合 \(\bigcap_{\alpha \in J}^{-1} f_{\alpha \gamma}(x_{\alpha})\) 有一个最大元素，等于 \(\inf_{\alpha \in J} (M_{\alpha \gamma}(x_{\alpha}))\) ；此外， \(G_{\gamma}\) 与 \(\bigcap_{\alpha \in J}^{-1} f_{\alpha \gamma}(x_{\alpha})\) 的交集是（由假设非空的）所有 \(y_{\gamma} \in G_{\gamma}\) 使得
\end{enumerate}

\[
y _ {\gamma} \leqslant \inf _ {\alpha \in J} \left(\mathbf {M} _ {\alpha \gamma} (x _ {\alpha})\right)
\]

的集合（利用条件(i)）。

\begin{enumerate}
\def\labelenumi{(\alph{enumi})}
\setcounter{enumi}{1}
\tightlist
\item
  让 \(J\) 是 \(I\) 的任意子集。一个族 \(x_{J} = (x_{\alpha})_{\alpha \in J}\) ，其中 \(x_{\alpha} \in G_{\alpha}\) 对于所有 \(\alpha \in J\) ，称为一致的，如果 \(x_{J}\) 的每个有限子族都是一致的。若 \(J \neq I\) 并且如果 \(\beta \in I - J\) ，证明存在 \(x_{\beta} \in G_{\beta}\) ，使得该族 \(x_{J \cup \{\beta\}} = (x_{\alpha})_{\alpha \in J \cup \{\beta\}}\) 是一致的。（利用(a)和条件(iv)，证明对每个有限子集 \(F\) （ \(J\) 的有限子集），如果 \(\gamma\) 是 \(F \cup \{\beta\}\) 的一个上界，那么 \(f_{\beta \gamma} \left( G_{\gamma} \cap \bigcap_{\alpha \in F} f_{\alpha \gamma}(x_{\alpha}) \right)\) 是（非空的）所有 \(y_{\beta} \in G_{\beta}\) ——它们 \(\leqslant f_{\beta \gamma} \left( \inf_{\alpha \in F} (M_{\alpha \gamma}(x_{\alpha})) \right)\) ——所组成的集合。利用 \(E_{\beta}\) 赋予相反序后是诺特集这一事实，接下来证明存在一个有限子集 \(F_0\) （ \(J\) 的有限子集）以及一个上界 \(\gamma_0\) （ \(F_0 \cup \{\beta\}\) 的上界），使得对于每个有限子集 \(F\) （ \(J\) 的有限子集）以及每个上界 \(\gamma\) （ \(F \cup \{\beta\}\) 的上界），我们拥有
\end{enumerate}

\[
f _ {\beta \gamma} \left( \right.\inf _ {\alpha \in \mathbf {F}} \left(\mathrm{M} _ {\alpha j} (x _ {\alpha})\right) \geqslant f _ {\beta \gamma_ {0}} \left(\inf _ {\alpha \in \mathbf {F} _ {0}} \left(\mathrm{M} _ {\alpha \gamma_ {0}} (x _ {\alpha})\right)\right).
\]

然后证明每个元素 \(x_{\beta} \in \mathbf{F}_{\beta}\) （它 \(\leqslant f_{\beta \gamma_0} \left( \inf_{\alpha \in \mathbf{F}_0} (\mathbf{M}_{\alpha \gamma_0}(x_\alpha)) \right)\) ）都满足所需条件。

\begin{enumerate}
\def\labelenumi{(\alph{enumi})}
\setcounter{enumi}{2}
\tightlist
\item
  最后，通过证明存在一个指标集为整个 I 的一致族来完成证明。（对一致族 \(x_{\mathbf{J}}\) 的集合通过关系“ \(x_{\mathbf{J}}\) 是 \(x_{\mathbf{K}}\) 的子族''进行排序，并应用（b）和佐恩引理。）
\end{enumerate}

\begin{enumerate}
\def\labelenumi{\arabic{enumi}.}
\setcounter{enumi}{5}
\item
  设 I 是一个有向集， \((\mathbf{J}_{\lambda})_{\lambda \in \mathbf{L}}\) 是满足习题1各条件的 I 的子集族。设 \((\mathbf{E}_{\alpha}, f_{\beta \alpha})\) 是由 I 索引的集合直接系统， \(\mathbf{E} = \lim \mathbf{E}_{\alpha}\) ；对每个 \(\lambda \in \mathbf{L}\) ，令 \(\mathbf{F}_{\lambda}\) 为把 \((\mathbf{E}_{\alpha}, f_{\beta \alpha})\) 的指标集限制到 \(J_{\lambda}\) 后所得直接系统的直接极限。每当 \(\lambda \leqslant \mu\) 时，令 \(g_{\mu \lambda}\) 为从 \(\mathbf{F}_{\lambda}\) 到 \(\mathbf{F}_{\mu}\) 的典范映射（第 6 号）。证明 \((\mathbf{F}_{\lambda}, g_{\mu \lambda})\) 是关于 L 的集合直接系统，并定义一个从 E 到 \(\mathbf{F} = \lim \mathbf{F}_{\lambda}\) 的典范双射。
\item
  设 I 是一个有向集，并令 \((\mathbf{E}_{\alpha}, f_{\beta\alpha})\) 是相对于I的一个集合直接系统。对于每个 \(\alpha \in I\) ，让 \(f_{\alpha}: E_{\alpha} \to E = \lim E_{\alpha}\) 为典范映射。在每一个 \(E_{\alpha}\) 中，让 \(R_{\alpha}\) 为等价关系 \(f_{\alpha}(x) = f_{\alpha}(y)\) .
\end{enumerate}

证明，每当 \(\alpha \leqslant \beta\) 时，映射 \(f_{\theta \alpha}\) 与等价关系 \(\mathbf{R}_{\alpha}\) 和 \(\mathbf{R}_{\beta}\) 相容。设 \(\mathbf{E}_{\alpha}^{\prime} = \mathbf{E}_{\alpha} / \mathbf{R}_{\alpha}\) ，并令 \(f_{\beta \alpha}^{\prime}\) 为从 \(\mathbf{E}_{\alpha}^{\prime}\) 到 \(\mathbf{E}_{\beta}^{\prime}\) 中的映射（过渡到商集后由 \(f_{\beta \alpha}\) 诱导）。证明 \(f_{\beta \alpha}^{\prime}\) 是单射，且 \((\mathbf{E}_{\alpha}^{\prime}, f_{\beta \alpha}^{\prime})\) 是一个集合直接系统，并定义一个从 \(\mathbf{E}\) 到 \(\lim_{\alpha \to 0} \mathbf{E}_{\alpha}^{\prime}\) 的典范双射。

\begin{enumerate}
\def\labelenumi{\arabic{enumi}.}
\setcounter{enumi}{7}
\item
  设 \((\mathbf{E}_{\alpha}, f_{\beta\alpha})\) 和 \((\mathbf{F}_{\alpha}, g_{\beta\alpha})\) 是两个集合直接系统，均由同一个有向集 I 索引。对每个 \(\alpha \in I\) ，令 \(u_{\alpha}\) 为从 \(E_{\alpha}\) 到 \(F_{\alpha}\) 的映射，并使各 \(u_{\alpha}\) 构成一个映射直接系统。设 \(G_{\alpha} \subset E_{\alpha} \times F_{\alpha}\) 是 \(u_{\alpha}\) 的图。证明 \((\mathbf{G}_{\alpha})\) 是各 \(E_{\alpha} \times F_{\alpha}\) 的一个子集直接系统，并且其直接极限可典范地等同于 \(u = \lim u_{\alpha}\) 的图。
\item
  设 I 为任意预序集， \((\mathrm{E}_{\alpha})_{\alpha\in\mathbf{I}}\) 为以 I 为指标集的集合族。对每一对指标 \((\alpha,\beta)\) （满足 \(\alpha\leqslant\beta\) ），令 \(f_{\beta\alpha}\) 为从 \(E_{\alpha}\) 到 \(E_{\beta}\) 的映射，并假设这些映射满足条件 \((\mathrm{LI}_{\mathbf{I}})\) 和 \((\mathrm{LI}_{\mathbf{II}})\) 。设 G 为族 \((\mathrm{E}_{\alpha})\) 的和所组成的集合，并且（采用第 5 号的记号）设 \(R\{x,y\}\) 为 G 的两个元素 x, y 之间的关系“ \(\lambda(x)=\alpha\leqslant\lambda(y)=\beta\) 且 \(y=f_{\beta\alpha}(x)\) ”。设 \(R'\) 是 G 上的等价关系，其图是包含 R 之图的所有等价关系的图中最小的一个（第二章，§6，习题 10）。集合 \(E=G/R'\) 称为族 \((\mathrm{E}_{\alpha})\) 关于映射族 \((f_{\beta\alpha})\) 的直接极限，记作 \(E=\lim_{\rightarrow}E_{\alpha}\) 。当指标集 I 有向时，证明此定义与第 5 号给出的定义一致。在一般情形下，从 G 到 E 的典范映射在 \(E_{\alpha}\) 上的限制，称为从 \(E_{\alpha}\) 到 E 的典范映射，记作 \(f_{\alpha}\) 。假设对每个 \(\alpha\in I\) 给定一个映射 \(u_{\alpha}\) （从 \(E_{\alpha}\) 到 F 中），使得 \(u_{\beta}\circ f_{\beta\alpha}=u_{\alpha}\) （每当 \(\alpha\leqslant\beta\) ）；证明存在唯一的从 E 到 F 的映射 u，使得 \(u=u_{\alpha}\circ f_{\alpha}\) （对每个 \(\alpha\in I\) ）.
\end{enumerate}
